\section{Problem Formulation}
\label{sec:problem}

\subsection{Events, Habits, and Decisions}

Let a session be an ordered stream of events $E=(e_1,\ldots,e_T)$.
Each event is represented as
\begin{equation}
e_t=(\tau_t,a_t,v_t,o_t,m_t)
\end{equation}
Here, $\tau_t$ is time, $a_t$ is the actor, $v_t$ is an action verb, $o_t$ is an object or entity, and $m_t$ contains domain metadata.
Structured tool calls are mapped deterministically; free-form language can be parsed into the same schema by an LLM.
The benchmark uses pre-structured events so event parsing is not part of the reported decision accuracy.

A habit $h$ represents a recurring behavioral regularity:
\begin{equation}
h=(g_h,c_h,k_h,q_h,f_h,\rho_h,z_h,V_h)
\end{equation}
Here, $g_h$ is a trigger action--entity pattern, $c_h$ is trigger context, $k_h$ records consequence statistics over observed typical outcomes, $q_h$ is confidence, $f_h$ and $\rho_h$ encode frequency and recency, $z_h$ is lifecycle state, and $V_h$ is a set of generalized variants.
Habit types include workflow, temporal, emotional, and associative patterns.

At each turn, the system observes the current event and memory state $M_t$, and outputs
\begin{equation}
\hat{y}_t = \pi(e_{\leq t},M_t) \in \{0,1\}
\end{equation}
Here, $1$ denotes a structured Push decision.
When $\hat{y}_t=1$, the decision also identifies its trigger, supporting habit, candidate score, and control trace.
Natural-language wording is a downstream host-system concern and does not change $\hat{y}_t$.

\subsection{Evaluation Objective}

Each evaluation turn has a ground-truth label $y_t$ indicating whether an intervention is useful at that stage.
We report precision, recall, F1, false positives (FP), and Push count in the main text; complete traces additionally retain accuracy and true-positive, false-negative, and true-negative counts.
F1 captures the balance between recovering useful interventions and avoiding false ones, while FP and Push count directly report false and total intervention frequency; they are offline proxies for potential intervention burden rather than measurements of perceived burden, satisfaction, trust, or helpfulness.
No single metric fully represents user utility: the relative cost of one missed intervention and one unnecessary interruption is deployment dependent.
We therefore treat threshold and cooldown as an operating policy rather than tune them solely to maximize test F1.

